\ysrc{543}我们来分析 Moore 复形上的结构。存在退化映射 $s_i\colon M_n\to M_{n+1}$（$0\le i\le n$）。拓扑直观提示：把退化单形「去掉」不应影响同调。为落实这一点，我们定义 $DM_*\subset M_*$，使 $DM_{n+1}$ 由退化映射 $s_0,\dots,s_n$ 的像张成。

\medskip\noindent\textbf{练习 A.4.4.} 验证 $DM_*$ 是 $M_*$ 的子复形，即它在微分 $\partial$ 之下保持不变。进一步，证明映射 $M_*\to(M/DM)_*$ 是拟同构。

基本上，我们想转而考虑链复形 $(M/DM)_*$。不过，使用它的另一种正规化更为方便，我们称之为\emph{正规化 Moore 复形}。定义 $NM_n$ 为\emph{所有} $i<n$ 的面映射 $d_i$ 的核（但 $i=n$ 不在其中）。于是 $(-1)^nd_n$ 定义了微分 $NM_n\to NM_{n-1}$。

\medskip\noindent\textbf{练习 A.4.5.} 验证 $NM_*$ 是链复形。

\begin{rem}{A.4.6}\label{thm:a6-07-1}
求和映射
\[
NM_n\oplus DM_n\to M_n
\]
是同构。因此，$NM_*$ 同构地映到 $(M/DM)_*$。
\end{rem}

从单纯 $R$-模到 $R$-模链复形的函子（我们将在定理 A.4.1 中采用）是 $\underline M\mapsto NM_*$。Dold--Kan 等价证明中的关键一步，是证明
\[
\underline{M}[n]\cong\bigoplus_{[n]\twoheadrightarrow[k]}NM_k
\]
关于 $\underline M$ 具有函子性。这里指标集取遍单纯范畴中满射的映射 $[n]\to[k]$（即保序映射）。这告诉我们如何定义逆函子：给定链复形 $M_*$，我们用一个单纯 $R$-模
\[
\underline{M}[n]:=\bigoplus_{[n]\twoheadrightarrow[k]}M_k
\]
来定义它。验证这样做确实给出范畴等价，其细节留给参考文献。

\subsection*{A.5.\quad 单纯交换环}

\emph{单纯交换环}是从 $\Delta^{\mathrm{op}}$ 到交换环范畴的函子 $\mathcal R$。由于我们将频繁处理这类对象，我们采用一个略微更经济的记号。单纯交换环用花体字母表示，并简记 $\mathcal R_n:=\mathcal R[n]$。经典交换环则用罗马字母如 $R$ 表示。

于是，一个单纯交换环可以具体地由一族交换环 $\mathcal R_n$（$n\ge0$）给出，连同它们之间由 $\Delta^{\mathrm{op}}$ 中态射索引的映射；这些映射又可由满足单纯恒等式的面映射 $d_i\colon\mathcal R_{n+1}\to\mathcal R_n$ 与退化映射 $s_i\colon\mathcal R_n\to\mathcal R_{n+1}$ 指定。这是对拓扑交换环的奇异单形上存在的结构所作的公理化。交换环的范畴记作 $\mathsf{CR}$，单纯交换环的范畴记作 $\mathsf{SCR}$。

\begin{exam}{A.5.1}\label{thm:a6-07-2}
设 $R$ 是交换环。则我们可以定义一个单纯交换环 $\underline R$，使 $\underline R[n]=R$，且所有面映射与退化映射都是恒等映射。直观上，$\underline R$ 对应于装备离散拓扑的拓扑交换环 $R$。
\end{exam}

\ysrc{544}\subsubsection*{A.5.1.\quad 同伦群}

若 $\mathcal R$ 是单纯交换环，则其同伦群 $\pi_*(\mathcal R)$ 构成分次交换环。$\pi_*(\mathcal R)$ 的底层群与 $\mathcal R$ 的底层单纯阿贝尔群的底层群一致，因此可由 Moore 复形
\[
\cdots\to\mathcal R_n\xrightarrow{\partial_n}\mathcal R_{n-1}\xrightarrow{\partial_{n-1}}\cdots\xrightarrow{\partial_0}\mathcal R_0
\]
计算。

\begin{exam}{A.5.2}\label{thm:a6-07-3}
对常值单纯交换环 $\underline R$，Moore 复形读作
\[
\cdots\xrightarrow{\mathrm{Id}}R\xrightarrow{0}R\xrightarrow{\mathrm{Id}}R\xrightarrow{0}R.
\]
于是 $\pi_0(\underline R)=R$，且 $i>0$ 时 $\pi_i(\underline R)=0$，这与例 A.5.1 中表达的直观一致。
\end{exam}

\begin{rem}{A.5.3}\label{thm:a6-07-4}
实际上，增广理想 $\pi_{>0}(\mathcal R)$ 在 $\pi_*(\mathcal R)$ 中总带有\emph{除幂结构}。我们不解释这意味着什么，也不解释它为何存在，只指出：对 $\Q$-代数，除幂结构总存在且唯一；而在正特征情形，它是一种相当特殊的结构。
\end{rem}

\subsubsection*{A.5.2.\quad 经典截断}

回忆 §2.3 中我们断言了伴随
\begin{equation}\tag{A.5.1}\label{eq:a6-07-1}
\begin{tikzcd}[column sep=5.4em, row sep=1.2em]
\mathsf{CR} \arrow[r, bend right=22, "{R\mapsto\underline R}"'] & \mathsf{SCR} \arrow[l, bend right=22, "{\pi_0(\mathcal R)\leftarrow\mathcal R}"']
\end{tikzcd}
\end{equation}
为论证这一点，我们需要说明
\[
\Hom_{\mathsf{SCR}}(\mathcal R,\underline R)\cong\Hom_{\mathsf{CR}}(\pi_0(\mathcal R),R).
\]
一个态射 $f\in\Hom_{\mathsf{SCR}}(\mathcal R,\underline R)$ 相当于一族映射 $f_n\colon\mathcal R_n\to R$，它们与面映射和退化映射相容。由于 $\underline R$ 上所有面映射与退化映射都是恒等映射，$f_n(r)$ 必然可以这样计算：先与 $n+1$ 个互不相同的映射 $\mathcal R_n\to\mathcal R_0$ 中的任一个复合，再作用 $f_0\colon\mathcal R_0\to R$。要使这一计算良定义，我们看到 $f_0$ 必须在 $d_1-d_0$ 的像上取零，从而经由 $\pi_0(\mathcal R)\to R$ 分解。剩下要说明的是，任何这样的 $f_0$ 确实诱导一个良定义的单纯交换环态射；这由下面的练习完成。

\medskip\noindent\textbf{练习 A.5.4.} 证明：由 $n+1$ 个映射 $[0]\to[n]$ 诱导的 $n+1$ 个映射 $\mathcal R_n\to\mathcal R_0$，与商映射 $\mathcal R_0\twoheadrightarrow\pi_0(\mathcal R)$ 复合后全都相同。

\subsubsection*{A.5.3.\quad 在单纯集上的富化}

由于这一构造的重要性，我们把单纯交换环在单纯集上的富化明确写出来，尽管它只是 §A.2.4 中描述的一个特例。给定集合 $S$ 与交换环 $R$，我们可以造出交换环 $R^{\otimes S}$。进而，对单纯集 $\underline S$ 与单纯交换环 $\mathcal R$，我们通过
\[
(\mathcal R^{\underline{\otimes}\underline S})_n:=(\mathcal R_n)^{\underline{\otimes}\underline S[n]}
\]
定义 $\mathcal R^{\underline{\otimes}\underline S}$。最后，我们定义一个单纯集 $\underline{\Hom}_{\mathsf{SCR}}(\mathcal R,\mathcal R')$，其中
\[
\underline{\Hom}_{\mathsf{SCR}}(\mathcal R,\mathcal R')[n]:=\Hom_{\mathsf{SCR}}(\mathcal R^{\otimes\Delta^n},\mathcal R'),
\]
并带有自然的面映射与退化映射。

\ysrc{545}\medskip\noindent\textbf{练习 A.5.5.} 对交换环 $R$ 与 $R'$，计算 $\underline{\Hom}_{\mathsf{SCR}}(\underline R,\underline R')$ 的同伦群。（答案：$\pi_0=\Hom_{\mathsf{CR}}(R,R')$，而 $i>0$ 时 $\pi_i$ 消没。）这为下述看法提供了依据：$\mathsf{CR}$ 以全忠实的方式嵌入 $\mathsf{SCR}$，此处「嵌入」是在单纯范畴的意义下说的。

\subsection*{A.6.\quad 单纯消解}

我们现在讨论一个关键构造：其中单纯交换环的态射被一个与投射消解或内射消解的形成相类似的过程「消解」。

\subsubsection*{A.6.1.\quad 导出函子}

让我们回忆同调代数中导出函子的范式，读者对此大概并不陌生。若 $F$ 是阿贝尔范畴 $\mathsf C$ 上的右正合函子，则为了在 $M\in\mathsf C$ 上定义 $F$ 的「高阶导出函子」，我们要找一个「自由」消解
\[
\cdots\to P_n\to P_{n-1}\to\cdots\to P_0\to M\to0,
\]
然后把 $F$ 作用到复形 $\cdots\to P_n\to P_{n-1}\to\cdots\to P_0$ 上；我们把这个复形看作 $M$ 的「替换」。至于「自由」这个形容词应当是什么意思，我们有意含糊其辞。

\begin{exam}{A.6.1}\label{thm:a6-07-5}
设 $\mathsf C$ 是 $R$-模范畴，$N$ 是 $R$-模，$F=N\otimes_R(-)$。则 $F$ 的高阶导出函子为 $\Tor_i^R(N,-)$。
\end{exam}

类似地，若 $F$ 是左正合函子，则改用\emph{内射}消解施行一个类比的过程。

\begin{exam}{A.6.2}\label{thm:a6-07-6}
设 $\mathsf C$ 是 $R$-模范畴，$N$ 是 $R$-模，$F=\Hom_R(N,-)$。则 $F$ 的高阶导出函子为 $\Ext^i_R(N,-)$。
\end{exam}

单纯交换环之于交换环，将扮演 $R$-模的链复形之于 $R$-模所扮演的角色。特别地，我们将看到：用某种「自由消解」替换交换环，允许我们在单纯交换环的范畴上建立一套「导出函子」理论。我们将在余切复形这一具体例子中说明这一点。

回忆：对态射 $f\colon A\to B$，我们有 Kähler 微分构成的 $B$-模 $\Omega^1_f$；当 $f$ 光滑时，它「性质良好」。对一般的 $f$，我们将用一个「光滑」的单纯 $A$-代数去「消解」$B$，以便定义余切复形 $\mathbf L_f$。在阿贝尔范畴中，这意味着用逐层具有良好性质（例如「自由」）的对象复形来替换 $B$；但这对 $A$-代数范畴并不适用，因为后者远非阿贝尔范畴。取而代之的是 Quillen 发展出的\emph{同伦代数}理论。其正确的框架是 Quillen 的\emph{模型范畴}，但若要完整解释，会带来相当大的离题，所以我们先采用一种权宜的表述，稍后在 §A.7 中提示一般理论。

\begin{defn}{A.6.3}\label{thm:a6-07-7}
单纯交换环的态射 $f\colon\mathcal R\to\mathcal R'$ 称为\emph{弱等价}，若它诱导同伦群的同构 $\pi_*(f)\colon\pi_*(\mathcal R)\xrightarrow{\sim}\pi_*(\mathcal R')$。
\end{defn}

这一概念是链复形的\emph{拟同构}概念在单纯交换环中的类比。事实上，我们也可以对单纯 $R$-模作类似定义；在 Dold--Kan 定理下，它会转化为通常的拟同构概念。

\ysrc{546}\subsubsection*{A.6.2.\quad 自由单纯代数}

设 $\mathcal A$ 是单纯交换环。一个\emph{单纯 $\mathcal A$-代数}是带有单纯交换环同态 $\mathcal A\to\mathcal R$ 的单纯交换环 $\mathcal R$。

\begin{defn}{A.6.4}\label{thm:a6-07-8}
\emph{自由单纯 $\mathcal A$-代数} $\mathcal R$ 是如下形式的单纯 $\mathcal A$-代数：
\begin{itemize}
\item 存在一族集合 $X_n$，使 $\mathcal R_n=\mathcal A[X_n]$；
\item 退化映射把 $s_j(X_n)$ 送入 $X_{n+1}$ 中，即 $s_j(X_n)\subset X_{n+1}$。（注意：面映射\emph{没有}任何条件！）
\end{itemize}
\end{defn}

下面的引理说明，自由单纯 $\mathcal A$-代数具有很好的「向外映射」性质，类似于模的自由消解所具有的性质。

\begin{lem}{A.6.5}\label{thm:a6-07-9}
考虑单纯代数的交换图
\[
\begin{tikzcd}[column sep=3.2em, row sep=2.4em]
\mathcal A \arrow[r] \arrow[d] & \mathcal B \arrow[d, "{\sim}\ \Phi"] \\
\mathcal R \arrow[r] \arrow[ur, dashed] & \mathcal C
\end{tikzcd}
\]
其中 $\Phi$ 逐层满射且是弱等价，而 $\mathcal R$ 是自由单纯 $\mathcal A$-代数。则存在提升 $\mathcal R\to\mathcal B$，使图表交换。
\end{lem}

\begin{proof}
见 [GS07, (5.4)]。关于模型范畴框架下更一般的陈述，见 [DS95, §4]。
\end{proof}

\begin{prop}{A.6.6}（消解的存在性）\label{thm:a6-07-10}
设 $f\colon A\to B$ 是（经典）交换环之间的映射。则存在自由单纯 $A$-代数 $\mathcal B$ 以及图表
\[
\begin{tikzcd}[column sep=3.6em, row sep=1.8em]
& \mathcal B \arrow[dr, "{\sim}"] & \\
\underline A \arrow[ur] \arrow[rr, "f"'] & & \underline B
\end{tikzcd}
\]
我们提醒：符号 $\xrightarrow{\sim}$ 表示弱等价，$\underline A$（相应地 $\underline B$）表示 $A$（相应地 $B$）上的常值单纯交换环。此外，我们甚至可以安排该图表的构造关于 $B$ 具有函子性。
\end{prop}

\begin{proof}[证明梗概]
我们将给出一个「典范」消解，它关于 $B$ 具有函子性。注意：对任何 $A$-代数 $R$，都有典范的代数同态
\begin{equation}\tag{A.6.1}\label{eq:a6-07-2}
A[R]\to R,
\end{equation}
它把 $[r]$ 送到 $r$。这一构造给出两个映射 $A[A[B]]\to A[B]$：
\begin{itemize}
\item 以 $R=B$ 套用 (A.6.1)，得到映射 $A[B]\to B$，再取 $A[-]$。例如，该映射把 $[a[b]]\mapsto[ab]$。
\item 以 $R=A[B]$ 套用 (A.6.1)。例如，该映射把 $[a[b]]\mapsto a[b]$。
\end{itemize}
此外还有一个映射 $A[B]\to A[A[B]]$，把 $[b]$ 送到 $[[b]]$；它是上述两个映射的截面。

细察此情形的组合结构，我们得到一个单纯对象
\begin{equation}\tag{A.6.2}\label{eq:a6-07-3}
\cdots A[A[B]]\rightrightarrows A[B],
\end{equation}
其中先前定义的映射就是最初的面映射与退化映射。显然这个消解逐层是 $A$ 上的多项式代数，并且可以验证它是自由的。为什么它是 $B$ 的消解呢？

\ysrc{547}这来自一个「额外退化」论证，它是一个更一般模式的一部分。若有一个单纯集
\[
\underline X=\cdots X_1\rightrightarrows X_0
\]
带有增广 $d_0\colon X_0\to X_{-1}$ 与「额外退化」$s_{-1}\colon X_{n-1}\to X_n$，它们满足单纯恒等式的自然延拓，则映射 $\underline X\to X_{-1}$（其中 $X_{-1}$ 视为常值单纯集）是弱等价。

回到 (A.6.2)，存在额外退化 $B\to A[B]$（$b\mapsto[b]$）与 $s_{-1}\colon A[B]\to A[A[B]]$（$a[b]\mapsto[a[b]]$），等等。由此可以证明：从 (A.6.2) 到 $B$ 的映射是弱等价。
\end{proof}

\begin{rem}{A.6.7}\label{thm:a6-07-11}
上述证明中的构造是一个更一般构造的特例；我们给出一般陈述，因为觉得它能澄清当前情形。设 $T$ 是范畴 $\mathsf C$ 上的一个\emph{单子}，即数据：一个自函子 $T\colon\mathsf C\to\mathsf C$，以及满足各种融贯条件的自然变换 $\eta\colon\mathrm{Id}_{\mathsf C}\to T$ 与 $\mu\colon T\circ T\to T$，细节见 [Mac71, §VI]。我们指出，单子在实践中常常作为伴随函子的复合而出现。（当前关注的典型例子是：$\mathsf C$ 是集合的范畴，$T$ 是把集合 $S$ 送到 $A[S]$ 的底层集合的函子。这里 $T$ 可以看作来自一个自由–遗忘伴随。）

单子 $T$ 上的一个\emph{代数}由数据 $x\in\mathsf C$ 以及满足融贯条件的态射 $a\colon Tx\to x$ 组成。（在上面的典型例子中，$T$ 上的代数范畴等价于 $A$-代数的范畴。）给定 $T$ 上的代数 $x$，可以造出一个称为 \emph{bar 构造}的单纯集 $B(T,x)_\bullet$，其中 $B(T,x)_n=T^{\circ(n+1)}x$。由形式的组合分析，bar 构造带有额外退化。因此，在很一般的情形下，增广 $B(T,x)_\bullet\to x$ 将是弱等价。
\end{rem}

\subsubsection*{A.6.3.\quad 唯一性}

若 $f\colon\mathcal A\to\mathcal B$ 是单纯交换环的态射，则对于图表
\[
\begin{tikzcd}[column sep=3.6em, row sep=1.8em]
& \mathcal B' \arrow[dr, "{\sim}"] & \\
\mathcal A \arrow[ur] \arrow[rr, "f"'] & & \mathcal B
\end{tikzcd}
\]
（其中 $\mathcal B'$ 是自由单纯 $\mathcal A$-代数），我们称之为 $\mathcal B$ 作为单纯 $\mathcal A$-代数的\emph{自由消解}。下面讨论这些消解「有多唯一」。在同调代数的经典情形中，投射消解在同伦意义下唯一。类似地，在某种意义下，单纯交换 $\mathcal A$-代数的自由消解也「在同伦意义下唯一」。表述这类陈述的最佳语言是模型范畴理论；完整解释它会使我们离题太远，但我们可以说明它在这一具体情形中特化为什么。

我们记 $\mathcal A[X]:=\mathcal A\otimes_{\underline{\Z}}\underline{\Z}[X]$，这是 $\mathcal A$ 上的一元多项式代数；其中张量积逐层取（本段中形成的所有张量积都如此；它们与 §5.2.3 中将要讨论的导出张量积是「同一个」东西）。我们称 $\mathcal A[X_0,X_1]:=\mathcal A[X]\otimes_{\mathcal A}\mathcal A[X]$。它表示「同伦余积」；特别地，给定两个 $\mathcal A$-代数同态 $f,g\colon\mathcal A[X]\to\mathcal B$，我们得到一个 $\mathcal A$-代数同态 $f\otimes g\colon\mathcal A[X_0,X_1]\to\mathcal B$。还存在乘积态射 $\mu\colon\mathcal A[X_0,X_1]\to\mathcal A[X]$。设 $\mathcal A[X_0,X_1,Y]$ 是 $\mathcal A[X_0,X_1]\xrightarrow{\mu}\mathcal A[X]$ 的单纯消解。

\ysrc{548}于是我们说 $f$ 与 $g$ \emph{同伦}，若 $f\otimes g$ 可延拓为一个交换图
\[
\begin{tikzcd}[column sep=4.2em, row sep=2.6em]
\mathcal A[X_0,X_1] \arrow[r] \arrow[dr, "{f\otimes g}"'] & \mathcal A[X_0,X_1,Y] \arrow[d, dashed] \\
& \mathcal B
\end{tikzcd}
\]
这样的图表称为 $f$ 与 $g$ 之间的一个\emph{同伦}。

我们说两个单纯消解
\[
\begin{tikzcd}[column sep=2.6em, row sep=1.6em]
& \mathcal B' \arrow[dr, "{\sim}"] & \\
\mathcal A \arrow[ur] \arrow[rr, "f"'] & & \mathcal B
\end{tikzcd}
\qquad
\begin{tikzcd}[column sep=2.6em, row sep=1.6em]
& \mathcal B'' \arrow[dr, "{\sim}"] & \\
\mathcal A \arrow[ur] \arrow[rr, "f"'] & & \mathcal B
\end{tikzcd}
\]
\emph{同伦}，若存在映射 $\alpha\colon\mathcal B'\to\mathcal B''$ 与 $\beta\colon\mathcal B''\to\mathcal B'$，使相应的图表交换，并且 $\alpha\circ\beta$ 与 $\beta\circ\alpha$ 都同伦于恒等。可以利用引理 A.6.5 证明：任意两个单纯消解都同伦。

\subsection*{A.7.\quad 模型范畴一瞥}

让我们提示这些程序背后更一般的框架，即 Quillen 的\emph{模型范畴}理论。我们的讨论将非常简短；一本受欢迎的入门参考书是 [DS95]。

模型范畴的目的是把同伦论纳入范畴论。基本想法是：在实践中，人们常常想引入一类「弱等价」映射，它们未必是同构。例如，拓扑空间的范畴中有通常的同伦等价概念，链复形的范畴中有拟同构概念。模型范畴就是在这样的情形下做同伦论构造的框架。

更形式地说，模型范畴是装备了若干被指定的态射类、并满足各种性质的范畴。其中最重要的是\emph{弱等价}，它们指认那些「同伦等价」的映射。弱等价正是我们想要「取逆」的态射。熟悉的例子有：
\begin{itemize}
\item 在 $R$-模链复形的范畴中，弱等价应当是拟同构。对它们取逆，就得到\emph{$R$-模的导出范畴}。
\item 在单纯交换环的范畴中，弱等价应当按定义 A.6.3 取定。
\end{itemize}
模型范畴中的其他数据是称为\emph{余纤维化}与\emph{纤维化}的态射类，它们大致是投射消解与内射消解的推广。这些态射类上带有各种公理：例如，它们在复合与收缩之下保持封闭。一个对象称为\emph{纤维性的}，若它到终对象的映射是纤维化；一个对象称为\emph{余纤维性的}，若它从始对象出发的映射是余纤维化。

\begin{exam}{A.7.1}\label{thm:a6-07-12}
在单纯集上的标准模型结构中，§A.2.3 提到的 \emph{Kan 复形}就是纤维性对象。Kan 复形由提升条件刻画，而拓扑空间的奇异单纯集正具有这些提升条件。
\end{exam}

\ysrc{549}特别地，对任何拓扑空间 $Y$，$\Sing Y$ 是 Kan 复形。这就是为什么在把单纯集的同伦群定义为从球面出发的映射时，必须把目标替换为一个 Kan 复形。

另一方面，在拓扑空间上的标准模型范畴结构中，纤维化是 Serre 纤维化；特别地，所有拓扑空间都是纤维性的。这就是为什么在定义拓扑空间的同伦群时，\emph{不}需要替换目标。

模型范畴的公理蕴涵：每个对象都有一条从某个余纤维性对象出发的弱等价，也有一条到某个纤维性对象去的弱等价。我们把这样的映射看作由余纤维性对象或纤维性对象给出的「消解」。我们知道，阿贝尔范畴上的导出函子正是用这样的消解定义的。模型范畴结构在不一定是阿贝尔的范畴中提供了更一般的「消解」概念，因此允许构造「非阿贝尔导出函子」。

我们关于单纯交换环的具体构造，是模型范畴中一般构造的例子。特别地，我们所说的「自由单纯 $\mathcal A$-代数 $\mathcal R$」，意思是 $\mathcal R$ 在标准模型结构中是余纤维性的 $\mathcal A$-代数；与此同时，所有单纯交换环都是纤维性的。

\subsection*{A.8.\quad 余切复形}

现在我们可以定义态射 $f\colon A\to B$ 的\emph{余切复形}了；它由 Quillen 发明，并由 Illusie 在 [Ill71, Ill72] 中发展。设 $\mathcal B\xrightarrow{\sim}B$ 是 $B$ 作为自由单纯 $A$-代数的单纯消解。再令 $\mathbf L_{\mathcal B/A}$ 为逐层取 Kähler 微分所得的单纯 $B$-模：
\[
\mathbf L_{\mathcal B/A}[n]:=\Omega^1_{\mathcal B_n/A}.
\]
最后，我们定义\emph{余切复形}为
\[
\mathbf L_f:=\mathbf L_{B/A}:=\mathbf L_{\mathcal B/A}\otimes_{\mathcal B}B.
\]

\medskip\noindent\textbf{练习 A.8.1.} 证明：若 $\mathcal B$ 与 $\mathcal B'$ 是 $B$ 的两个自由单纯消解，则 $\mathbf L_{\mathcal B/A}\otimes_{\mathcal B'}B$ 同伦于 $\mathbf L_{\mathcal B'/A}\otimes_{\mathcal B}B$。因此，$\mathbf L_{B/A}$ 在同伦意义下良定义。特别地，它给出 $B$-模导出范畴中一个良定义的对象。

\begin{rem}{A.8.2}（导出推广）\label{thm:a6-07-13}
更好的是：若 $\mathcal A$ 是动态环，$\mathcal B$ 是动态 $\mathcal A$-代数，则 $\mathbf L_{\mathcal B/\mathcal A}$ 可以构造为一个动态 $\mathcal B$-模——见 [Lur04, §3.2]。
\end{rem}

当态射 $f$ 具有良好的有限性性质时，我们期待 $\mathbf L_{B/A}$ 也具有合理的有限性性质。然而，命题 A.6.6 的证明中用来构造 $\mathcal B$ 的典范消解极其巨大。因此，我们希望知道存在「更小」的消解。我们下面就来考察这个问题。

\subsection*{A.9.\quad 经济的消解}

我们将按照 [Iye07, §4] 证明具有「良好」有限性性质的消解的存在性。

\begin{prop}{A.9.1}\label{thm:a6-07-14}
设 $A$ 是诺特交换环，$f\colon A\to B$ 是有限型态射。则存在 $B$ 的单纯 $A$-代数消解，它由自由单纯 $A$-代数 $A[X]$ 给出，且每个 $X_n$ 都是有限集。
\end{prop}

\begin{cor}{A.9.2}\label{thm:a6-07-15}
设 $A$ 是诺特交换环，$f\colon A\to B$ 是有限型态射。则 $\mathbf L_{B/A}$ 可由有限自由 $B$-模的复形表示。
\end{cor}

\ysrc{550}\begin{rem}{A.9.3}\label{thm:a6-07-16}
一般而言，$\mathbf L_{B/A}$ 不能由完美复形表示；也就是说，它通常在无穷多个次数上有同调群。例如，若 $A$ 是域且 $\mathbf L_{B/A}$ 的 Tor-幅度有限，则实际上 $\mathbf L_{B/A}$ 的 Tor-幅度至多为 $2$，并且 $B$ 在 $A$ 上是 LCI。这是 Avramov 的结果 [Avr99]，它解决了 Quillen 的一个猜想。
\end{rem}

作为动机，我们回忆如何用自由 $A$-模的复形为有限型 $A$-模 $M$ 构造「高效」消解。我们可以构造一列有限自由 $A$-模的复形 $\{F_i\}$，其同调在次数不超过 $i$ 的范围内逼近 $M$。对 $F_0$，我们任取一个从自由模出发的满射 $A^n\twoheadrightarrow M$。然后取 $I:=\ker(F_0\twoheadrightarrow M)$ 的生成元，这诱导一个像为 $I$ 的映射 $A^m\xrightarrow{\partial}A^n$，而它到 $M$ 的映射在 $H_0$ 上诱导同构。于是我们可以取 $F_0=[A^m\xrightarrow{\partial}A^n]$。接着，我们由 $F_{i-1}$ 归纳地构造 $F_i$：在 $F_{i-1}$ 中为 $H_i(F_{i-1})$ 的生成元选取代表元，再在次数 $i+1$ 处添加自由直和项，使这些生成元成为边界。

要对单纯交换环施行类似的构造，我们需要一种在单纯交换环中「消去闭链」的办法。下面的引理实现了这一点。

\begin{lem}{A.9.4}\label{thm:a6-07-17}
给定单纯交换环 $\mathcal A$ 与同伦类 $[z]\in\pi_i(\mathcal A)$，存在自由单纯 $\mathcal A$-代数 $\mathcal A'$，使得：
\begin{itemize}
\item 对每个 $n$，$\mathcal A'_n$ 是有限生成自由 $\mathcal A_n$-代数，且 $n\le i$ 时 $\mathcal A'_n=\mathcal A_n$；
\item 映射 $\mathcal A\to\mathcal A'$ 在 $\pi_n$（$n<i$）上诱导同构，而在 $n=i$ 时给出正合列
\[
0\to\pi_0(\mathcal A)\cdot[z]\to\pi_i(\mathcal A)\to\pi_i(\mathcal A')\to0.
\]
\end{itemize}
\end{lem}

确实，假设引理 A.9.4 成立，我们便能像命题 A.9.1 那样构造消解：取一个满射 $\mathcal A[x_1,\dots,x_d]\twoheadrightarrow B$，然后反复套用引理 A.9.4，构造出一列自由单纯 $\mathcal A$-代数，它们在所有有限次数上都具有「正确的」有限性性质与同伦群。取它们的余极限，就得到所要的消解。

现在来证明引理 A.9.4。我们想「添加」一个边界为 $[z]$ 的变量。然而由于单纯恒等式，变量「添加」的过程必然相当复杂。为此，我们模仿向拓扑空间 $X$ 粘胞腔以消去次数 $i-1$ 的同调类时，奇异单形上发生的事情。若在次数 $i$ 上粘一个胞腔，就会在所有更高维中产生额外的退化单形。由于这个原因，该构造要复杂得多。

\begin{proof}[引理 A.9.4 的证明]
我们定义
\[
X_n:=\{x_t\mid t\colon[n]\twoheadrightarrow[i+1]\in\Delta\},
\]
并取 $\mathcal A'_n=\mathcal A_n[X_n]$，面映射与退化映射如下定义。
\begin{itemize}
\item 有 $s_j(x_t)=x_{t\circ\sigma_j}$。
\item 注意 $X_{i+1}=\{x_{\mathrm{Id}}\}$。我们令 $d_0(x_{\mathrm{Id}})=z$，并对 $j>0$ 令 $d_j(x_{\mathrm{Id}})=0$。对 $n>i+1$，若 $t\circ\delta_j$ 满射，则定义 $d_j(x_t)=x_{t\circ\delta_j}$；否则 $t\circ\delta_j$ 会经由一个面映射 $d_{j'}\colon[i]\to[i+1]$ 分解，于是我们定义它，使得
\ysrc{551}下图交换
\[
\begin{tikzcd}[column sep=3.4em, row sep=2.4em]
{}[n] \arrow[r, "t"] \arrow[d, "{\delta_j}"'] & {}[i+1] \\
{}[n-1] \arrow[r] & {}[i] \arrow[u, "{\delta_{j'}}"']
\end{tikzcd}
\]
\end{itemize}
可以验证这样得到的 $\mathcal A'_n$ 是满足引理 A.9.4 结论的单纯交换环（更多细节见 [Matc]）。
\end{proof}

\begin{rem}{A.9.5}\label{thm:a6-07-18}
用更概念化的方式陈述该构造如下（沿 [Matc]）。对单纯 $A$-代数 $\mathcal A$，$\pi_i(\mathcal A)$ 中的一个类可由单纯集的映射 $\partial\Delta^{i+1}\to\mathcal A$ 表示（这里用了注 A.3.3：$\mathcal A$ 是纤维性的）。这诱导单纯交换环之间的映射 $A[\partial\Delta^{i+1}]\to\mathcal A$，其中 $A[\partial\Delta^{i+1}]$ 是单纯集 $\partial\Delta^{i+1}$ 上的自由单纯 $A$-代数。于是我们可以构造
\[
\mathcal A':=\mathcal A\otimes_{A[\partial\Delta^{i+1}]}A[\Delta^{i+1}],
\]
并断言它满足引理 A.9.4 的结论。第一个项目符号由构造直接满足。为计算它对同伦群的影响，我们注意 $\pi_*(\mathcal A)\xrightarrow{\sim}\pi_*(A[\Delta^{i+1}])$，因为 $\Delta^{i+1}$ 可缩；而 $A[\partial\Delta^{i+1}]$ 的最低正次数同伦群出现在次数 $i$，且按构造映到 $\pi_i(\mathcal A)$ 中的 $[z]$。然后用 Tor 谱序列 [Qui70, (5.2)] 作结论：
\[
\begin{aligned}
&\Tor^{\pi_*(A[\partial\Delta^{i+1}])}_p\bigl(\pi_*(\mathcal A),\pi_*(A[\Delta^{i+1}])\bigr)_q\\
&\qquad\Longrightarrow\ \pi_{p+q}\bigl(\mathcal A\otimes_{A[\partial\Delta^{i+1}]}A[\Delta^{i+1}]\bigr).
\end{aligned}
\]
\end{rem}

\subsection*{A.10.\quad 导出张量积}

设 $\mathcal A\to\mathcal B$ 与 $\mathcal A\to\mathcal B'$ 是单纯交换环之间的两个映射。设 $\mathcal A\hookrightarrow\mathcal P\xrightarrow{\sim}\mathcal B$ 是 $\mathcal B$ 作为 $\mathcal A$-代数的自由消解，$\mathcal A\hookrightarrow\mathcal P'\xrightarrow{\sim}\mathcal B'$ 是 $\mathcal B'$ 作为 $\mathcal A$-代数的自由消解。\emph{导出张量积}「$\mathcal B\overset{\mathbf L}{\otimes}_{\mathcal A}\mathcal B'$」由单纯交换环 $\mathcal P\otimes_{\mathcal A}\mathcal P'$ 表示，其中
\[
(\mathcal P\otimes_{\mathcal A}\mathcal P')_n=\mathcal P_n\otimes_{\mathcal A_n}\mathcal P'_n.
\]
虽然它并不唯一（因为它依赖于消解的选取），但它在弱等价意义下唯一。事实上，只需消解其中一项即可，例如，单纯交换环 $\mathcal P\otimes_{\mathcal A}\mathcal B'$ 同样表示该导出张量积。确实，存在显然的映射
\[
\mathcal P\otimes_{\mathcal A}\mathcal P'\to\mathcal P\otimes_{\mathcal A}\mathcal B',
\]
并且它是自由 $\mathcal B'$-代数之间的弱等价。

若 $A$ 是经典环，$\mathcal P$ 是自由单纯 $A$-代数，则 $\mathcal P$ 的底层单纯 $A$-模在 Dold--Kan 对应下对应于一个自由 $A$-模的复形。特别地，Dold--Kan 对应把自由单纯消解送到链复形中人们熟悉的自由消解概念。特别地，我们看到：若 $\mathcal A\xrightarrow{\sim}A$、$\mathcal B\xrightarrow{\sim}B$、$\mathcal B'\xrightarrow{\sim}B'$ 全都经典，则
\[
\pi_i(\mathcal B\overset{\mathbf L}{\otimes}_{\mathcal A}\mathcal B')\cong\Tor_i^A(B,B').
\]
这有如下推论。

\begin{cor}{A.10.1}\label{thm:a6-07-19}
设 $A\to B$ 与 $A\to B'$ 是经典交换环之间的映射，且对一切 $i>0$ 有 $\Tor_i^A(B,B')=0$。则
\[
B\overset{\mathbf L}{\otimes}_A B'\to\pi_0(B\overset{\mathbf L}{\otimes}_A B')\cong B\otimes_A B'
\]
是弱等价。
\end{cor}

\ysrc{552}\subsection*{A.11.\quad 余切复形的性质}

现在我们来确立余切复形的一些性质；这些性质对 Kähler 微分的模来说是很熟悉的，至少光滑情形下如此。作为参考，我们推荐 [Qui]。

\begin{prop}{A.11.1}\label{thm:a6-07-20}
若 $B,B'$ 是 $A$-代数，且对 $i>0$ 有 $\Tor_i^A(B,B')=0$，则
\begin{enumerate}[(1)]
\item $\mathbf L_{B\underline{\otimes}_AB'/B'}\cong\mathbf L_{B/A}\otimes_A B'$。
\item $\mathbf L_{B\underline{\otimes}_AB'/A}\cong(\mathbf L_{B/A}\otimes_A B')\oplus(B\otimes_A\mathbf L_{B'/A})$。
\end{enumerate}
\end{prop}

\begin{rem}{A.11.2}\label{thm:a6-07-21}
一般而言，下述结论成立：
\begin{enumerate}[(1)]
\item $\mathbf L_{B\overset{\mathbf L}{\underline{\otimes}}_AB'/B'}\cong\mathbf L_{B/A}\otimes_A B'$。
\item $\mathbf L_{B\overset{\mathbf L}{\underline{\otimes}}_AB'/A}\cong(\mathbf L_{B/A}\otimes_A B')\oplus(B\otimes_A\mathbf L_{B'/A})$。
\end{enumerate}
Tor 消没假设正是用来保证 $B\overset{\mathbf L}{\otimes}_AB'\xrightarrow{\sim}B\otimes_A B'$。
\end{rem}

\begin{proof}
两个论证都可从 Kähler 微分关于张量积的标准性质形式地推出，所用的观察是：
\begin{itemize}
\item 若 $\mathcal P$ 是自由单纯 $\mathcal A$-代数，$\mathcal B$ 是任一单纯 $\mathcal A$-代数，则 $\mathcal P\otimes_{\mathcal A}\mathcal B$ 是自由单纯 $\mathcal B$-代数。
\end{itemize}
\end{proof}

\begin{prop}{A.11.3}\label{thm:a6-07-22}
若 $A\to B\to C$ 是态射的复合，则在 $C$-模的导出范畴中存在正合三角
\[
\mathbf L_{B/A}\otimes_B C\to\mathbf L_{C/A}\to\mathbf L_{C/B}.
\]
\end{prop}

\begin{proof}
取自由消解 $A\hookrightarrow\mathcal B\xrightarrow{\sim}B$。这赋予 $C$ 一个 $\mathcal B$-代数结构，于是我们可以再取自由消解 $\mathcal B\hookrightarrow\mathcal C\xrightarrow{\sim}C$。这就给出图表
\[
\begin{tikzcd}[column sep=3.4em, row sep=1.8em]
& \mathcal C \arrow[ddr, "{\sim}"] & \\
& \mathcal B \arrow[u, hook] \arrow[d, "{\sim}"] & \\
A \arrow[ur, hook] \arrow[r] & B \arrow[r] & C
\end{tikzcd}
\]
然后我们作张量积
\[
\begin{tikzcd}[column sep=3.4em, row sep=1.8em]
& \mathcal C \arrow[ddr, "{\sim}"] \arrow[dr] & \\
& \mathcal B \arrow[u, hook] \arrow[d, "{\sim}"] & B\otimes_{\mathcal B}\mathcal C \arrow[d] \\
A \arrow[ur, hook] \arrow[r] & B \arrow[r] \arrow[ur] & C
\end{tikzcd}
\]
对光滑态射的 Kähler 微分正合列 [Sta22, Tag 02K4] 逐层应用，我们得到正合三角
\[
\mathbf L_{\mathcal B/\mathcal A}\otimes_{\mathcal B}\mathcal C\to\mathbf L_{\mathcal C/\mathcal A}\to\mathbf L_{\mathcal C/\mathcal B}.
\]
再作用 $-\otimes_{\mathcal C}C$，便得到正合三角。下面验证各项确实如所断言。

\ysrc{553}\begin{itemize}
\item $\mathbf L_{B/A}\otimes_B\mathcal C\otimes_{\mathcal C}C\cong\mathbf L_{B/A}\otimes_B C\cong\mathbf L_{B/A}\otimes_B\mathcal C$。
\item $\mathbf L_{\mathcal C/A}\otimes_{\mathcal C}C\cong\mathbf L_{\mathcal C/A}$。
\item 有 $\mathbf L_{\mathcal C/\mathcal B}\otimes_{\mathcal B}B\cong\mathbf L_{B\otimes_{\mathcal B}\mathcal C/B}$。由于 $B\otimes_{\mathcal B}\mathcal C$ 是 $C$ 作为 $B$-代数的自由消解，在 $C$ 上以 $\mathcal C$ 作张量便给出 $\mathbf L_{C/B}$。
\end{itemize}
\end{proof}

\begin{cor}{A.11.4}（局部化的余切复形）\label{thm:a6-07-23}
设 $A$ 是交换环。
\begin{enumerate}[(1)]
\item 若 $S$ 是 $A$ 中的乘法系，则 $\mathbf L_{S^{-1}A/A}\cong0$。
\item 设 $A\to B$ 是环同态。若 $S$ 是 $A$ 中的乘法系，$T$ 是 $B$ 中包含 $S$ 的像的乘法系，则
\[
\mathbf L_{T^{-1}B/S^{-1}A}\cong\mathbf L_{B/A}\otimes_B T^{-1}B.
\]
\end{enumerate}
\end{cor}

\begin{proof}
(1) 乘法诱导 $S^{-1}A\otimes_A S^{-1}A\xrightarrow{\sim}S^{-1}A$。把它代入命题 A.11.1(ii)，得到
\[
\mathbf L_{S^{-1}A/A}\cong\mathbf L_{S^{-1}A\otimes_AS^{-1}A/A}\cong\mathbf L_{S^{-1}A/A}\oplus\mathbf L_{S^{-1}A/A},
\]
经由对角映射。这迫使 $\mathbf L_{S^{-1}A/A}\cong0$。

(2) 由命题 A.11.1(i)，有 $\mathbf L_{S^{-1}B/S^{-1}A}\cong\mathbf L_{B/A}\otimes_B S^{-1}B$。套用命题 A.11.3，得到正合三角
\[
\mathbf L_{S^{-1}B/S^{-1}A}\otimes_{S^{-1}B}T^{-1}B\to\mathbf L_{T^{-1}B/S^{-1}A}\to\mathbf L_{T^{-1}B/S^{-1}B}.
\]
最右项按 (i) 消没\tn{原文作「according to (i)」，应为 (1)（原书 p.553）。}，由此即得结论。
\end{proof}

\subsection*{A.12.\quad 余切复形的一些计算}

现在我们在一些特殊情形中辨认余切复形。单纯消解几乎总是过于笨重，无法手工计算，因此我们转而必须巧妙地利用上面解释的形式性质。

\begin{prop}{A.12.1}\label{thm:a6-07-24}
设 $A$ 是诺特交换环，$f\colon A\to B$ 是有限型态射。则
\begin{enumerate}[(1)]
\item $f$ 平展，当且仅当 $\mathbf L_f\cong0$。
\item $f$ 光滑，当且仅当 $\mathbf L_f\cong\Omega^1_f$ 且为有限投射。
\end{enumerate}
\end{prop}

\begin{proof}[一个方向的证明]
目前我们只证明正向。反向将等到发展了 $\mathbf L_f$ 与形变理论之间的联系之后再确立。

设 $f$ 是平展的。则乘法映射 $B\otimes_A B\to B$ 是局部化，故推论 A.11.4 蕴涵 $\mathbf L_{B/B\otimes_A B}\cong0$。由命题 A.11.1，有 $\mathbf L_{B\otimes_A B/B}\cong\mathbf L_{B/A}\otimes_A B$。于是对 $B\xrightarrow{\mathrm{Id}_B\otimes1_B}B\otimes_A B\to B$ 套用命题 A.11.3，得到 $\mathbf L_{B/A}\cong0$。

接下来设 $f$ 是光滑的。光滑映射的局部结构定理说：在 $\Spec B$ 上局部地，$f$ 可以分解为一个仿射空间上的平展映射：
\[
\begin{tikzcd}[column sep=3.6em, row sep=2em]
\Spec B \arrow[r] \arrow[dr, "f"'] & \mathbb A^n_A \arrow[d] \\
& \Spec A
\end{tikzcd}
\]
在环的层面上，这意味着经过适当的局部化后，我们有分解 $A\to P\to B$，其中 $P$ 是 $A$ 上的多项式环，$P\to B$ 平展。

\ysrc{554}此时，常值单纯 $A$-代数的态射 $\underline A\to P$ 已经是自由的，因此 $\mathbf L_{P/A}\cong\Omega_{P/A}$ 是有限自由模。由命题 A.11.3，我们有正合三角 $\mathbf L_{P/A}\otimes_P B\to\mathbf L_{B/A}\to\mathbf L_{B/P}$，其中 $\mathbf L_{P/A}\otimes_P B\cong\Omega_{B/A}$、$\mathbf L_{B/P}=0$，这就完成了证明。
\end{proof}

回忆：交换环 $A$ 中的\emph{正则序列}是一个有限序列 $r_1,\dots,r_n\in A$，使得每个 $r_i$ 都是模 $A/(r_1,\dots,r_{i-1})$ 的非零因子。

\begin{prop}{A.12.2}\label{thm:a6-07-25}
若 $B=A/I$，且 $I$ 由正则序列生成，则
\[
\mathbf L_{B/A}\xrightarrow{\sim}I/I^2[1].
\]
\end{prop}

\begin{proof}
先处理 $A=\Z[x]$、$B=\Z$、$I=(x)$ 的情形。此时我们有截面 $\Z\to\Z[x]$。对序列 $\Z\to\Z[x]\to\Z$ 套用命题 A.11.3，得到正合三角
\[
\mathbf L_{\Z[x]/\Z}\otimes_{\Z[x]}\Z\to\mathbf L_{\Z/\Z}\to\mathbf L_{\Z/\Z[x]}.
\]
这说明 $\mathbf L_{\Z/\Z[x]}\cong\mathbf L_{\Z[x]/\Z}[1]\otimes_{\Z[x]}\Z\cong(x)/(x^2)$。

现在考虑命题的更一般情形。由归纳法，只需处理 $I=(r)$ 为主理想的情形。注意 $r$ 的选取诱导映射 $\Z[x]\to A$（$x\mapsto r$），而该映射诱导 $A/(r)\cong\Z\otimes_{\Z[x]}A$。我们断言：$r$ 是非零因子，当且仅当对一切 $i>0$ 有 $\Tor_i^{\Z[x]}(\Z,A)=0$；换言之（由推论 A.10.1），当且仅当 $\Z\overset{\mathbf L}{\otimes}_{\Z[x]}A\xrightarrow{\sim}A/(r)$。该断言是同调代数中的初等计算，所用消解为 $\Z$ 在 $\Z[x]$ 上的消解 $\Z[x]\xrightarrow{x}\Z[x]$。于是由命题 A.11.1，我们得到
\[
\mathbf L_{B/A}\cong\mathbf L_{\Z/\Z[x]}\otimes_{\Z[x]}A\cong(x)/(x^2)[1]\otimes_{\Z[x]}A\cong(r)/(r^2)[1].
\]
\end{proof}

\subsection*{A.13.\quad 整体化}

我们讨论的理论可以整体化到概形，进而整体化到叠。例如，设 $Y\to X$ 是概形的态射。则我们这样定义 $\mathbf L_{Y/X}$：把 $\mathcal O_Y$ 作为单纯 $\mathcal O_X$-代数作自由替换，然后逐层取 $\Omega$，等等。传统上，以这种方式发展该理论（而不是试图从仿射开子空间「粘接」余切复形）很重要，原因在于范畴技术的不足：一方面，余切复形在链复形层面上并不唯一；另一方面，导出范畴不适合粘接。

然而，用现代的语言，确实可以按所希望的方式仿射局部地粘接余切复形：把余切复形看作动态模，并利用动态拟凝聚层的范畴满足 Zariski 下降这一事实。

\subsection*{A.14.\quad André--Quillen 同调}

设 $A\to B$ 是交换环之间的映射。我们已经把余切复形 $\mathbf L_{B/A}$ 定义为一个在同伦意义下的单纯 $B$-模。

\begin{defn}{A.14.1}\label{thm:a6-07-26}
设 $M$ 是 $B$-模。我们定义 \emph{André--Quillen 同调群}
\[
D_i(B/A;M):=H_i(\mathbf L_{B/A}\otimes_B M),
\]
以及 \emph{André--Quillen 上同调群}
\[
D^i(B/A;M):=H^i(\Hom_B(\mathbf L_{B/A},M)).
\]
\end{defn}
